\chapter{Marco Teórico}

\section{Datos faltantes en ciencia de datos}

En el ámbito de la ciencia de datos y el aprendizaje automático, se denomina
valor faltante a la ausencia de un valor en una o más variables de un registro
del conjunto de datos. Esta situación es común en datasets del mundo real y
puede originarse por diversas causas: fallos en sensores o dispositivos de
medición, errores humanos durante la recolección, pérdida de información en
la transmisión de datos, o simplemente porque ciertas variables no aplican
para determinados registros \parencite{little2019}.

La presencia de valores faltantes representa un problema fundamental en el
preprocesamiento de datos, dado que la mayoría de los algoritmos de
aprendizaje automático requieren conjuntos de datos completos para su
funcionamiento. Ignorar los registros con valores faltantes mediante
eliminación (\textit{listwise deletion}) puede derivar en una pérdida
significativa de información y en la introducción de sesgos, especialmente
cuando los datos no se pierden de forma aleatoria \parencite{garcia2017}.
Por esta razón, la imputación de valores faltantes se ha consolidado como
una etapa estándar en cualquier pipeline de ciencia de datos.

\section{Mecanismos de ausencia de datos}

La naturaleza de los valores faltantes no es siempre la misma
\parencite{little2019}  propone una clasificación formal de los mecanismos
de ausencia en tres categorías, cuya comprensión es fundamental para
elegir una estrategia de imputación adecuada.

\subsection{MCAR: Missing Completely At Random}

Se dice que los datos son MCAR cuando la probabilidad de que un valor sea
faltante es completamente independiente de cualquier variable del dataset,
tanto de los valores observados como de los valores ausentes:

\begin{equation}
    P(\mathbf{M} \mid \mathbf{X}_{obs}, \mathbf{X}_{mis}) = P(\mathbf{M})
\end{equation}

donde $\mathbf{M}$ es la matriz de indicadores de valores faltantes,
$\mathbf{X}_{obs}$ representa los valores observados y $\mathbf{X}_{mis}$
los valores ausentes. Bajo este mecanismo, los datos faltantes pueden
considerarse una muestra aleatoria de todos los datos, lo que lo convierte
en el escenario más favorable para la imputación \parencite{little2019}.
En el presente trabajo se asume este mecanismo para el diseño del
experimento controlado.

\subsection{MAR: Missing At Random}

Bajo el mecanismo MAR, la probabilidad de ausencia depende de los valores
observados, pero no de los valores faltantes en sí mismos. Por ejemplo,
la probabilidad de que el pH de una muestra de agua no sea registrado
podría depender de la temperatura del agua, que sí fue medida. Este
mecanismo es más común en la práctica que MCAR, aunque su verificación
empírica es compleja \parencite{little2019}.

\subsection{MNAR: Missing Not At Random}

En el mecanismo MNAR, la probabilidad de ausencia depende del propio valor
faltante. Este es el escenario más difícil de tratar, ya que la ausencia
de información está relacionada con la información misma. Por ejemplo, si
los valores extremos de una variable tienden a no ser registrados, los
datos faltantes no son aleatorios y cualquier método de imputación
introducirá sesgo \parencite{garcia2017}.

\section{Imputación de datos faltantes}

La imputación es el proceso mediante el cual los valores faltantes de un
dataset son estimados y reemplazados por valores calculados a partir de la
información disponible. El objetivo no es recuperar los valores originales
con exactitud, sino preservar la estructura estadística del dataset de forma
que el modelo de aprendizaje automático entrenado sobre los datos imputados
sea lo más representativo posible \parencite{salgado2016}.

Existen múltiples estrategias de imputación, que van desde métodos simples
basados en estadísticos descriptivos hasta enfoques más sofisticados basados
en similitud entre instancias o modelos probabilísticos. A continuación se
describen los dos métodos evaluados en el presente trabajo.

\subsection{Imputación por media}

La imputación por media es el método más utilizado en la práctica debido
a su simplicidad y eficiencia computacional \parencite{garcia2017}. Consiste
en reemplazar cada valor faltante $x_{ij}$ del atributo $j$ con la media
aritmética $\bar{x}_j$ calculada sobre los valores disponibles de dicho
atributo:

\begin{equation}
    \hat{x}_{ij} = \bar{x}_j = \frac{1}{n_j}
    \sum_{\substack{i=1 \\ x_{ij} \neq \emptyset}}^{n} x_{ij}
\end{equation}

donde $n_j$ es el número de instancias que tienen el atributo $j$ completo.

Este método presenta ventajas claras en términos de velocidad y facilidad
de implementación. Sin embargo, su principal limitación es que al asignar
un único valor global a todos los registros faltantes de una columna, reduce
artificialmente la varianza del atributo y destruye las correlaciones
existentes entre variables \parencite{little2019}. Esto puede traducirse en
un dataset imputado que no refleja fielmente la distribución original de los
datos, afectando negativamente el rendimiento del modelo.

\subsection{Imputación por K-Nearest Neighbors (KNN)}

La imputación por KNN es un método basado en similitud que supera una de
las principales limitaciones de la imputación por media: la ignorancia de
las relaciones entre variables \parencite{troyanskaya2001}. En lugar de
utilizar un estadístico global, este método estima el valor faltante de
una instancia a partir de los valores del atributo correspondiente en sus
$k$ vecinos más cercanos, identificados mediante distancia euclidiana en
el espacio de atributos observados:

\begin{equation}
    \hat{x}_{ij} = \frac{1}{k} \sum_{l \in \mathcal{N}_k(i)} x_{lj}
\end{equation}

donde $\mathcal{N}_k(i)$ es el conjunto de los $k$ vecinos más cercanos de
la instancia $i$ entre aquellas que presentan el atributo $j$ completo, y
$x_{lj}$ es el valor del atributo $j$ en el vecino $l$.

La intuición detrás de este enfoque es que instancias similares en el
espacio de características tienden a tener valores similares en sus
atributos. Al aprovechar esta estructura local del dataset, la imputación
KNN preserva en mayor medida las correlaciones entre variables, lo que
puede resultar en una mejor calidad de los datos imputados y, por
consiguiente, en un mejor rendimiento del clasificador \parencite{salgado2016}.

La elección del parámetro $k$ influye en el resultado: valores pequeños
de $k$ son más sensibles al ruido local, mientras que valores grandes
tienden a suavizar demasiado las estimaciones. En la literatura se reporta
que $k = 5$ es un valor de partida ampliamente adoptado \parencite{salgado2016},
y es el valor utilizado en el presente experimento.

\section{Modelos de clasificación supervisada}

Un modelo de clasificación supervisada es un algoritmo que aprende a
asignar una etiqueta de clase a nuevas instancias a partir de un conjunto
de entrenamiento etiquetado \parencite{bishop2006}. En el contexto de este
trabajo, el modelo de clasificación opera sobre los datos previamente
imputados y su rendimiento es utilizado como indicador indirecto de la
calidad de la imputación aplicada.

\subsection{Regresión Logística}

La Regresión Logística es uno de los modelos de clasificación binaria más
utilizados como línea base (\textit{baseline}) en ciencia de datos, debido
a su interpretabilidad y su bajo costo computacional \parencite{bishop2006}.
El modelo estima la probabilidad de que una instancia pertenezca a la clase
positiva mediante la función sigmoide:

\begin{equation}
    P(y = 1 \mid \mathbf{x}) = \frac{1}{1 + e^{-(\mathbf{w}^T \mathbf{x} + b)}}
\end{equation}

donde $\mathbf{w}$ es el vector de pesos aprendido durante el entrenamiento
y $b$ es el sesgo (\textit{bias}). La clase predicha se obtiene aplicando
un umbral de decisión, típicamente de 0.5.

\section{Métricas de evaluación}

La evaluación del rendimiento de un clasificador requiere el uso de métricas
que reflejen con precisión su capacidad predictiva. En el presente trabajo
se emplean dos métricas complementarias.

\subsection{Accuracy}

El Accuracy mide la proporción de predicciones correctas sobre el total de
instancias evaluadas:

\begin{equation}
    \text{Accuracy} = \frac{TP + TN}{TP + TN + FP + FN}
\end{equation}

donde $TP$, $TN$, $FP$ y $FN$ representan los verdaderos positivos,
verdaderos negativos, falsos positivos y falsos negativos, respectivamente.
Si bien es una métrica intuitiva, puede resultar engañosa en presencia de
clases desbalanceadas \parencite{japkowicz2002}.

\subsection{F1-score}

El F1-score es la media armónica entre la precisión (\textit{precision}) y
la exhaustividad (\textit{recall}):

\begin{equation}
    F_1 = 2 \cdot \frac{\text{precision} \cdot \text{recall}}
    {\text{precision} + \text{recall}}
\end{equation}

Esta métrica es especialmente adecuada cuando existe desbalance de clases,
ya que penaliza los modelos que ignoran la clase minoritaria y predicen
sistemáticamente la clase mayoritaria \parencite{japkowicz2002}. Por esta
razón, se considera la métrica principal de comparación en el presente
experimento.
